\begin{figure*}[t]
\centering
\includegraphics[width=\textwidth]{figures/fig1_overview.png}
\caption{Readout-dependent kernel equivalence, identifiable
dimension, and recovery limits.
(a) Different quantum states and measured features can define
the same Gaussian kernel. A common input-independent transformation
preserves all pairwise kernel values when it induces
$f_\eta(x)=Qf_\theta(x)+\tau$ with $Q^\top M Q=M$.
The congruent triangles illustrate the isotropic case;
arbitrary unitaries need not preserve the specified kernel.
(b) Schematic kernel-equivalence classes in parameter space:
points on each colored curve share the same all-input kernel
and coefficient Gram matrix $G$, whereas different curves
represent different kernels. The curves illustrate the equivalence
relation, not the dimension or global topology of a particular
circuit's equivalence classes.
For a fixed data basis with $\{1,h_1,\ldots,h_A\}$ linearly
independent on the actual input domain, fixed $M\succ0$ and
$\gamma>0$, and analytic coefficients on an open connected
parameter box, Theorem~\ref{thm:gram} characterizes all-input
kernel equivalence, global coordinate inactivity, and generic
differential rank through $G$.
(c) In the four-qubit, one-layer Hubregtsen/local-Z example,
fixed positive diagonal metrics give rank two, and fixed dense
positive-definite metrics reach but cannot exceed rank three.
Any fixed affine feature lift with a fixed positive-definite
quadratic metric also has rank at most three.
Adding the two same-input-group products $f_0f_2$ and $f_1f_3$
gives rank four under a positive diagonal metric.
Solid and dashed arrows encode the numbers of identifiable
and locally null directions in a stage-adapted basis; they are
not spatial coordinate axes or matched directions across stages.
The products are computable from existing exact local means
in this example, without new ideal measurement information.
All ranks are generic; two products suffice without a minimality
claim. Finite-shot resolution, predictive improvement, and
execution savings do not follow from rank alone.}
\label{fig:overview}
\end{figure*}